\documentclass[exercice,a4]{thedocument}%
\usepackage[french]{babel}%
\usepackage{pstricks}%
\usepackage{scratch3}%
\usepackage{tabularx}%
\usepackage{array}%
\startdatakeys
    \source{22GENMATMEAG1}
    \BO{2026}
    \cycle{4}
    \competencesDuSocle{CA,RA}
    \objectifsApprentissage{C4E1ag,C4B1ed}
    \motsCles{Scratch,programme,aire,développement,équation}
\enddatakeys
\begin{document}%

\medskip

Dans cet exercice, $x$ est un nombre strictement supérieur à 3.

On s'intéresse aux deux figures géométriques dessinées ci-dessous :
\begin{itemize}
\item un rectangle dont les côtés ont pour longueurs $x-3$ et $x+7$.
\item un carré de côté $x$.
\end{itemize}

\begin{center}
\begin{pspicture}(-1,-0.5)(13,4.5)
\psframe(0,0)(5,2)\psframe(7,0)(11,4)
\uput[l](0,1){$x - 3$}\uput[d](2.5,0){$x + 7$}\uput[l](7,2){$x$}\uput[d](9,0){$x$}
\end{pspicture}
\end{center}

\begin{enumerate}
\item Quatre propositions sont écrites ci-dessous :

Recopier sur la copie celle qui correspond à l'aire du carré. On ne demande pas de justifier.

\begin{tabularx}{\linewidth}{|*{4}{>{\centering \arraybackslash}X|}}\hline
$4x$&$4+x$&$x^2$&$2x$\\\hline
\end{tabularx}

\item  Montrer que l'aire du rectangle est égale à $x^2 + 4x - 21$.

\item On a écrit le script ci-dessous dans Scratch.

On veut que ce programme renvoie l'aire du rectangle lorsque l'utilisateur a rentré une valeur de $x$ (strictement supérieure à 3).

Écrire sur la copie les contenus des trois cases vides des lignes 5, 6 et 7, en précisant les numéros de lignes qui correspondent à vos réponses.

\hspace*{2cm}
\begin{scratch}[num blocks]
\blockinit{Quand la touche\selectmenu{espace} est pressée}
\blocklook{demander \ovalnum {combien vaut x ?} et attendre}
\blockvariable{mettre \selectmenu{x} à \ovalsensing{réponse}}
\blockvariable{mettre \selectmenu{R} à \ovalvariable{x}*\ovalvariable{x}}
\blockvariable{ajouter \ovalnum{}*\ovalvariable{x} à \ovalvariable{R}}
\blockvariable{ajouter \ovalnum{} à \ovalvariable{R}}
\blocklook{dire \ovaloperator{regrouper \ovalnum{L'aire du rectangle est }et \ovalnum{}} pendant \ovalnum{2} secondes}
\end{scratch}\\

\item On a pressé la touche espace puis saisi le nombre 8. Que renvoie le programme ?

\item Quel nombre $x$ doit-on choisir pour que l'aire du rectangle soit égale à l'aire du carré ?

\emph{Toute trace de recherche, même non aboutie, sera prise en compte.}
\end{enumerate}


\vspace{0.5cm}

\showthesource
\end{document}
